\documentclass[10pt,a4paper]{article}
\usepackage[T1]{fontenc}
\usepackage[margin=22mm]{geometry}
\usepackage{lmodern,mathrsfs,amsmath,amssymb,amsthm,mathtools,booktabs,microtype}
\usepackage[hidelinks]{hyperref}
\newtheorem{theorem}{Theorem}
\newtheorem{lemma}[theorem]{Lemma}
\title{Prime-power boundary comparison and two-sided ES moment certificates}
\author{The Clankers}
\date{14 September 2026\quad---\quad Research preprint}
\begin{document}\maketitle
\begin{abstract}
A complete labelled ES divisor atlas has a prime-labelled truncated-polynomial
module. Its joint kernel and ordinary augmentation have a comparison whose
trace is precisely the original success count. Two singularity-safe canonical
moment problems bound that trace from below and above, retaining split-fibre
masses and all relation norms. Four moments determine 9530 of 9531 shell counts
in the specified finite test; six moments resolve the remaining case. The
construction does not establish a universal ES theorem. No priority claim is made.
\end{abstract}

\paragraph{Arithmetic input and exact ordered inverse.}
Fix a prime $p\equiv1\pmod4$, an integer $p/4<a<p/2$, and $R=4a-p$.
Write $a=\prod_iq_i^{e_i}$ with distinct increasing primes. An original state
is a pair $(c,\beta)$, where $c=E,M$ and $-e_i\le\beta_i\le e_i$.
Keep $u=\prod_iq_i^{e_i+\beta_i}$, $g_E=4u+1$, $g_M=u+a$ and
\[
 D=\frac R{\gcd(R,g_c)},\qquad x=D^{-1}\in(0,1/3]\cup\{1\}.
 \tag{1}
\]
The two positive rational triples are
\[
 E:\ \left(a,\frac{pa+a^2/u}R,\frac{pa+p^2u}R\right),\qquad
 M:\ \left(a,\frac{p(a+a^2/u)}R,\frac{p(a+u)}R\right).
 \tag{2}
\]
They satisfy $(RY-pa)(RZ-pa)=p^2a^2$, hence $4/p=1/a+1/Y+1/Z$.
Since $4a\equiv p\pmod R$ and $\gcd(pa,R)=1$, the two reduced denominators
of $Y,Z$ are both exactly $D$. Thus $D=1$ is equivalent to an integral hit.
These are the inherited Type I/II divisor coordinates \cite{ET}.
The ordered inverse is $u=a^2/(RY-pa)$ in E and $u=pa^2/(RY-pa)$ in M.
With $d_0=\gcd(a,u)$, retain $(h,r,s)=(d_0^2/u,u/d_0,a/d_0)$.
Prime valuations prove $a=hrs$, $u=hr^2$, $\gcd(r,s)=1$.
The integer quotients at a hit are $(pr+s)/R$ in E and $(r+s)/R$ in M.
Neither channel nor denominator order is identified with the other.

\begin{theorem}[The boundary comparison detects original successes]
For a finite labelled selection $\mathscr A$ of original states, let $\mathcal L$
be all primes dividing its moduli and set
\[
 d_{\alpha\ell}=\max\{v_\ell(R_\alpha)-v_\ell(g_\alpha),0\},\quad
 J_\alpha=\mathbb Q[z_\ell:\ell\in\mathcal L]/(z_\ell^{d_{\alpha\ell}+1})_\ell.
\]
Write $\mathcal J=\bigoplus_\alpha J_\alpha$ and $V=\bigoplus_\alpha\mathbb Q e_\alpha$.
The maps
\[
 s:V\xrightarrow{\sim}\bigcap_\ell\ker(z_\ell:\mathcal J\to\mathcal J),\quad
 e_\alpha\mapsto z^{d_\alpha}e_\alpha,\qquad
 \epsilon(F_\alpha)=F_\alpha(0)e_\alpha
\]
obey $\epsilon s=\operatorname{diag}(\mathbf1_{D_\alpha=1})$.
Its trace and rank are exactly the number $H$ of original canonical hits.
\end{theorem}
\begin{proof}
The monomials $z^b$, $0\le b_\ell\le d_{\alpha\ell}$, form a basis.
Multiplication by each $z_\ell$ has distinct nonzero images, so its kernel
has exactly the basis vectors with $b_\ell=d_{\alpha\ell}$. Their joint
kernel is the top-monomial line. Reading that coefficient inverts $s$.
Evaluation at zero preserves it precisely when every $d_{\alpha\ell}$ is
zero, equivalently $D_\alpha=1$.
\end{proof}
This is the explicit multi-prime application of the source's boundary-kernel
map $M/vM\simeq\ker(v:M/vf(M)\to M/vf(M))$, $[x]\mapsto[f(x)]$,
with $M=\mathbb Q[z]$, $f=z^d$ \cite[BR1--BR7]{SZ}.
Each failure still has a nonzero top line; forgetting its degree is not success.
A supported zero retains its state label, rather than becoming external absence.
For $d'\le d$, ordinary reduction kills the top line unless $d'=d$;
the reverse injection $F\mapsto z^{d-d'}F$ preserves it. These are different maps.

\newpage
\paragraph{Canonical lower and upper bounds on the comparison trace.}
Let $\mu=\sum_\alpha w_\alpha\delta_{x_\alpha}$ be any finite positive measure
on $(0,c]\cup\{1\}$, with $0<c<1$. Retain its full mass, not a normalized
probability measure. Put $m_j=\sum_\alpha w_\alpha x_\alpha^j$ and
$H=\mu(\{1\})$. For $d\ge0$ set
\[
 U_d=\min_{\deg P\le d,\,P(1)=1}\int P^2\,d\mu,\qquad
 L_d=\max_{\deg P\le d,\,P(1)=1}\frac1{1-c}\int(x-c)P^2\,d\mu.
 \tag{3}
\]
The upper problem is the classical Christoffel variational problem
\cite[\S2.1, (2.2)--(2.3)]{Las}. The finite singular kernels, signed lower
problem and original arithmetic return are calculated here.
For $0\le i,j<d$ define
\[
 \begin{aligned}
 (G_+)_{ij}&=m_{i+j+2}-2m_{i+j+1}+m_{i+j},&
 (b_+)_i&=m_{i+1}-m_i,\\
 (G_-)_{ij}&=-m_{i+j+3}+(c+2)m_{i+j+2}-(2c+1)m_{i+j+1}+cm_{i+j},&
 (b_-)_i&=m_{i+2}-(c+1)m_{i+1}+cm_i.
 \end{aligned}\tag{4}
\]
\begin{theorem}[Two-sided certificate, including singular source forms]
Both $G_\pm$ are positive semidefinite and $b_\pm\in\operatorname{im}G_\pm$.
For any solutions $G_\pm y_\pm=b_\pm$,
\[
 \boxed{L_d=\frac{m_1-cm_0+b_-^{\mathsf T}y_-}{1-c}
 \ \le\ H\ \le\ U_d=m_0-b_+^{\mathsf T}y_+.}\tag{5}
\]
All optimizers are $P_+=1-(x-1)y_+(x)+(x-1)\ker G_+$ and
$P_-=1+(x-1)y_-(x)+(x-1)\ker G_-$.
The bounds are nested and become exact at a finite degree.
If $H$ counts original states, $\lceil L_d\rceil=\lfloor U_d\rfloor$
certifies its exact value. In particular $L_d>0$ forces a witness;
$U_d<1$ certifies absence only in the stated finite domain.
\end{theorem}
\begin{proof}
Write $P=1+(x-1)q$. The two quadratic expressions are
\[
 m_0+2b_+^{\mathsf T}q+q^{\mathsf T}G_+q,
 \quad (m_1-cm_0+2b_-^{\mathsf T}q-q^{\mathsf T}G_-q)/(1-c).
\]
The Gram weights are $w(x-1)^2$ and $w(c-x)(x-1)^2$, respectively,
nonnegative on the declared support. Any kernel polynomial vanishes where
that weight is positive. The corresponding $b$ weight vanishes at every
other point, so $b$ annihilates the kernel and lies in the range.
Completing the squares proves the formulas, attainment and complete optimizer
fibres. Successful atoms contribute $H$; other atoms give nonnegative upper
and nonpositive lower contributions. Increasing degree enlarges the feasible
space. A polynomial equal to one at 1 and vanishing at every other support
point proves finite exactness.
\end{proof}
No generalized inverse is hidden: rational RREF gives one solution and a
complete rational nullspace. The degree-zero convention uses empty matrices.
For $i\le j$, orthogonality to all allowed corrections proves the exact secants
\[
 U_i-U_j=\int(P_{+,i}-P_{+,j})^2d\mu,\qquad
 L_j-L_i=\frac1{1-c}\int(c-x)(P_{-,i}-P_{-,j})^2d\mu.
 \tag{6}
\]
Both differences have the retained primitive $(P_i-P_j)/(x-1)$.
The lower weight vanishes at $x=c$; these original states remain present
elsewhere in the source. This instantiates the original-relation norm
calculation of PR29 \cite[\S3--4]{SZ}; it does not identify counting measure
with the theta-function norm.

\paragraph{Four integer moments and a finite stopping estimate.}
For ES, $c=1/3$ and $m_j=S_j/R^j$, where
$S_j=\sum_\alpha\gcd(R,g_\alpha)^j$. Define
$A_i=3S_{i+1}-RS_i$ for $i=0,1,2$.
The signed degree-one matrix is negative semidefinite exactly when
\[
 A_0\le0,\quad A_2\le0,\quad A_0A_2\ge A_1^2.\tag{7}
\]
Any positive direction has nonzero evaluation at 1, because polynomials
vanishing there have nonpositive signed norm. Normalizing it proves that
violation of (7) is equivalent to $L_1>0$.
The comparison $P=T_d(6x-1)/T_d(5)$ also proves
$0\le U_d-H\le m_0/T_d(5)^2$ and
$0\le H-L_d\le m_0/(2T_d(5)^2)$.
Thus $T_d(5)^2>3m_0/2$ gives a unique integer count. This is a degree bound,
not a claimed fast algorithm for the original moments.

\newpage
\paragraph{Split cohomology, its two quotients, and unchanged mass.}
For an actual budget split $e=f+b$, the map $(\xi,\eta,c)\mapsto(\xi+\eta,c)$
has complete fibre
\[
 \max(-f_i,\beta_i-b_i)\le\xi_i\le\min(f_i,\beta_i+b_i),\quad
 \eta=\beta-\xi,
\]
whose size is the product $n(\beta)$ of those interval lengths.
Each occurrence is given the original state's mass divided by $n(\beta)$.
For every polynomial $P$,
\[
 \sum_\alpha w_\alpha P(x_\alpha)
 =\sum_\omega\frac{w_{q\omega}}{n(\beta)}P(x_{q\omega}).\tag{8}
\]
Pullback duplicates values; its weighted adjoint averages them over each
fibre. Their composite is the identity. Their other composite is the
orthogonal projection, with kernel the explicit fibre-sum-zero relations.
They commute with the prime variables and both maps of Theorem 1. This proves
that the moments and canonical bounds retain the original count under splitting.

The previous residue cohomology has a different quotient. For matching pairs
$\mathscr P$, original successes $S$, and common channel/residues $I$, keep
\[
 \mathbb Q^S\xleftarrow{q_*}\mathbb Q^{\mathscr P}
 \xrightarrow{r_*}\mathbb Q^I\cong H^0_{\rm residue}.
\]
The averaged comparison $B_{z\alpha}=|q^{-1}(\alpha)\cap r^{-1}(z)|/n(\beta)$
preserves total mass; its norm is $B^{\mathsf T}G_I B$, not an assigned identity.
At $p=2521,a=636,R=23$, $f=(1,1,1),b=(1,0,0)$, the three dimensions are
$3,5,4$. In column order $(E,477),(M,8),(M,50562)$ and row order
$(E,13),(M,11),(M,21),(M,22)$,
\[
 B=\begin{pmatrix}1&0&0\\0&1/2&0\\0&0&1/2\\0&1/2&1/2\end{pmatrix},\quad
 G_I=\operatorname{diag}(2,2,2,3/2),\quad
 B^{\mathsf T}G_I B=\begin{pmatrix}2&0&0\\0&7/8&3/8\\0&3/8&7/8\end{pmatrix}.
 \tag{9}
\]
The nonzero cross terms remain, as required by the mixed-support source formula
BR8 \cite{SZ}. Neither cohomology quotient is identified with the other.

\paragraph{Exact certificates.}
For this shell the measure is $3\delta_1+87\delta_{1/23}$; the boundary module
is $\mathbb Q^3\oplus(\mathbb Q[z]/z^2)^{87}$. Its full socle has dimension
90 but the comparison trace is three. The polynomial $(23x-1)/22$ gives
$L_1=U_1=3$, using four original moments.
For the empty shell $p=241,a=64,R=15$, the measure is
$8\delta_{1/3}+12\delta_{1/5}+6\delta_{1/15}$ and
\[
 -12/425\le H\le180/463<1.
\]
This certifies that shell empty, not the prime false. Its 26 supported socle
lines still exist.

All 9531 first-half shells at the 54 primes $p\equiv1\pmod{12}$, $p\le1500$,
were enumerated: 412838 original states and 1920 hits in 685 shells.
Degree-one brackets give exact counts in 9530 shells; the remaining case
$p=1093,a=280,R=27$ has
\[
 [L_1,U_1]=[45982/7737,681/94],\quad
 [L_2,U_2]=[6,1132982/188317],
\]
so the exact count is six. Every occupied shell in this test has $L_1>0$;
that reverse implication is not proved for arbitrary shells. The standard-library
verifier uses exact integers and fractions, singular normal equations,
full prime-exponent boxes and negative controls. It computes moments by
finite enumeration; no universal positivity, faster complexity, analytic
Frobenius comparison or Lean certification is asserted.

\begingroup\small
\begin{thebibliography}{3}
\bibitem{SZ} The Clankers (2026). \emph{SplitZero source continuations}.
Zeta repository: PR28 \texttt{tau-split-integration/RESEARCH\_NOTE.md};
PR29 \texttt{tau-arithmetic-metric-transfer/RESEARCH\_NOTE.md}, at
\texttt{35b7aeaf\ldots2507fb}. Draft PR30
\texttt{tau-mixed-boundary-formal/RESEARCH\_NOTE.md}, BR1--BR8,
at \texttt{f57532b9\ldots3e97}. Full identifiers and locators accompany the sources.
\bibitem{Las} Lasserre, J.-B. (2022). \emph{A disintegration of the Christoffel
function}. arXiv:2203.16238v1, \S2.1, equations (2.2)--(2.3).
\bibitem{ET} Elsholtz, C., \& Tao, T. (2013). Counting the number of solutions
to the Erd\H{o}s--Straus equation on unit fractions.
\emph{Journal of the Australian Mathematical Society, 94}(1), 50--105.
\S2, Propositions 2.2 and 2.6.
\end{thebibliography}
\endgroup
\end{document}
